\section{Demostració de la convergència}

La part més complicada ja l'hem superada, ara només queda escriure el productori infinit com el d'una successió $(1+a_n)_{n\geq 0}$ i aplicar el la proposició \ref{prop}.

\begin{thm}
    El producte $C(f_1,\ldots,f_k)$ convergeix sota les condicions de la conjectura de Bateman i Horn.
\end{thm}
\begin{proof}
    Notem que
    \begin{align*}
    \left(1-\frac{1}{p}\right)^{-k}\left(1-\frac{\omega_f(p)}{p}\right)&=\left(1+\frac{k}{p}+\frac{k(k+1)}{2p^2}+O\left(\frac {1}{p^3}\right)\right)\left(1-\frac{\omega_f(p)}{p}\right)\\
    &=1+\frac{k-\omega_f(p)}{p}+\frac{B(p)}{p^2},
    \end{align*}
    on 
    \[
    B(p):=\frac{k(k+1)}{2}-k\omega_f(p)+O\left(\frac 1 p \right).
    \]
    A més, per $p$ prou gran,
    \[
    \lvert B(p)\rvert\leq \frac{k(k+1)}{2}+k\deg(f)+C
    \]
    on $C$ és una constant donada per $O\left(1/p \right)$. És a dir, $B(p)$ està uniformement fitada. Prenem 
    \[
    a_p:=\frac{k-\omega_f(p)}{p}+\frac{B(p)}{p^2}.
    \]
    Aleshores,
    \[
    \sum_{p \text{ primer}}a_p=\sum_{p\text{ primer}}\frac{k-\omega_f(p)}{p}+\sum_{p\text{ primer}}\frac{B(p)}{p^2}
    \]
    convergeix. D'altra banda, 
    \[
    \lvert a_p\rvert^2=O\left(\frac 1 {p^2}\right).
    \]
    Així doncs, per la proposició \ref{prop}, 
    \[
    \prod_{p\text{ primer}}(1+a_p)=\prod_{p\text{ primer}}\left(1-\frac{1}{p}\right)^{-k}\left(1-\frac{\omega_f(p)}{p}\right)
    \]
    convergeix.
\end{proof}